\documentclass[12pt]{article}
\usepackage[brazil]{babel}
\usepackage[utf8]{inputenc}
\usepackage[T1]{fontenc}
\usepackage{amsmath, amssymb}
\usepackage{geometry}
\usepackage{graphicx}
\usepackage{xcolor}
\usepackage{enumitem}

\definecolor{maincolor}{RGB}{0,100,50}
\usepackage{fancyhdr}
\definecolor{forestgreen}{RGB}{0,100,50}
\usepackage[most]{tcolorbox}

\newtcolorbox{block}[1]{%
  colback=green!5,
  colframe=maincolor,
  coltitle=white,
  title=#1,
  fonttitle=\bfseries,
  boxrule=1.5pt,
  arc=3mm,
  left=6pt,
  right=6pt,
  top=6pt,
  bottom=6pt
}

\geometry{margin=2.5cm}
\pagestyle{fancy}
\fancyhf{}

\renewcommand{\headrulewidth}{2pt}
\renewcommand{\footrulewidth}{2pt}

\renewcommand{\headrule}{\hbox to\headwidth{%
  \color{forestgreen}\leaders\hrule height \headrulewidth\hfill}}

\renewcommand{\footrule}{\hbox to\headwidth{%
  \color{forestgreen}\leaders\hrule height \footrulewidth\hfill}}

\setlength{\headheight}{32pt}

\fancyhead[L]{%
  % Certifique-se de que o arquivo de imagem esteja na mesma pasta
  \includegraphics[height=1.2cm]{Universidade_de_Rio_Verde.jpg}
}

\fancyhead[R]{%
  \raggedleft
  \textcolor{forestgreen}{%
    \small
    \textbf{Lista de Exercícios 2.}\\
    Me.~Sergio Souza Novak\\
    \textit{Lógica para Engenharia de Software}
  }
}

\fancyfoot[C]{%
  \textcolor{forestgreen}{\thepage}
}

\begin{document}

\vspace*{0.2cm}
\section*{Exercícios — Teoria dos Conjuntos}



\begin{enumerate}

\item Considere um conjunto arbitrário \(A\). Prove formalmente as seguintes afirmações envolvendo o conjunto vazio, utilizando as definições de inclusão, união e interseção:

\begin{enumerate}[label=\alph*)]
    \item \(\varnothing \subseteq A\).
    \item \(A \cap \varnothing = \varnothing\).
    \item \(A \cup \varnothing = A\).
    \item Se \(A \subseteq \varnothing\), então \(A = \varnothing\).
    \item \(\varnothing \cap \varnothing = \varnothing\).
    \item \(\varnothing \cup \varnothing = \varnothing\).
\end{enumerate}

\item Prove a propriedade reflexiva da inclusão:
\[
A \subseteq A.
\]

\item Mostre que a inclusão é antissimétrica. Ou seja, se:
\[
A \subseteq B \quad \text{e} \quad B \subseteq A,
\]
então \(A = B\).

\item Prove que a interseção é um subconjunto dos conjuntos originais:
\[
A \cap B \subseteq A.
\]

\item Prove que os conjuntos originais são subconjuntos da união:
\[
A \subseteq A \cup B.
\]

\item Utilizando a definição de pertinência, mostre que, se \(x \in A \cap B\), então \(x \in A\) e \(x \in B\).

\item Mostre que, se \(x \in A\), então \(x \in A \cup B\).

\item Prove a propriedade comutativa da interseção:
\[
A \cap B = B \cap A.
\]

\end{enumerate}

\end{document}